\section{Hypothetical Design Scenarios}
\label{sec:case_studies}

The following scenarios illustrate how the vocabulary could guide design and evaluation. They are not case studies with collected data, implemented controllers, or demonstrated benefits. Each requires a domain-specific model and comparison with established practice.

\subsection{Clinical decision support}
A decision-support system may have to allocate further assessment while accounting for uncertainty, clinician workload, and the consequences of delay. Epistemic pull could represent the expected decision value of an additional observation; selective protection could concern integrity of the patient record and validated model components; computational drag could track inference latency and processing cost; action efficiency could represent timely decision support; and sustainability could track review capacity across a sequence of cases.

Clinical test costs and burdens should not be conflated with computation costs: they belong in the explicitly defined action and resource model. Clinical eligibility, escalation rules, and patient safety requirements would be externally specified constraints. No treatment recommendation or clinical benefit follows from the force metaphor. Evaluation would require clinical oversight, relevant approvals, and evidence that the model improves appropriate outcomes over validated decision-support baselines.

\subsection{Energy management}
For a resource-limited infrastructure controller, a candidate state includes reserve capacity, predicted demand, equipment condition, and available computation. The proposed decomposition separates forecasting, preservation of critical components, computation expenditure, dispatch performance, and future supply capacity. Existing optimization methods can represent these quantities; UVIF would need to show that its decomposition improves monitoring, adaptation, or diagnosis.

A useful comparison would apply identical demand sequences, reserve models, actuator limits, and forecast information to the competing controllers. Outcomes would include unmet demand, constraint violations, computational burden, and reserve depletion over a declared horizon. A favourable average cost would not compensate for an unreported violation of a hard operational requirement.

\subsection{Autonomous transportation}
The navigation example in Section~\ref{sec:optimization_equilibrium} can be extended to sensor allocation and memory updates. A UVIF implementation would need an explicit transition model, a way to connect policy updates to vehicle controls, and a safety mechanism that respects physical braking and sensing limitations. Equilibrium-seeking dynamics alone do not ensure recovery from an obstacle or sensor failure.

Testing should separate normal operation, known disturbances, and conditions outside the assumed model. Task completion, reserve use, intervention frequency, and failures should be reported together. Comparators should include resource-aware constrained planning or control, not only a controller optimized for travel speed.

\subsection{Cybersecurity}
A monitoring system may allocate inspection effort between immediate threat detection and sustained availability. The components could represent information gathering, protection of keys and logs, inspection expenditure, response effectiveness, and future monitoring capacity. A threat model must state attacker capabilities and which assets are protected. No architectural label makes a system resistant to adversarial inputs.

Ablations could examine whether an explicit capacity model reduces overload under repeated alerts, while retaining detection performance. Both missed threats and disruptions to legitimate activity matter. Tuning the protective component only on attacks used for evaluation would invalidate the intended comparison.

\subsection{Organizational capacity}
A service organization may balance information gathering, retention of expertise, processing effort, timely service, and future staffing capacity. Here the central modelling issue is governance: whose workload, risk, and service quality count, and who may alter the constraints? A persistent organization can still impose unacceptable costs on its members or clients.

The framework can make these trade-offs explicit, but it cannot select legitimate priorities on its own. Evaluation would require stakeholder-defined outcomes and comparisons with existing capacity-planning methods. Institutional longevity is not an adequate sole measure of success.
